\documentclass[11pt]{article}
\usepackage{amsmath,amssymb,amsthm}
\usepackage[margin=25mm]{geometry}
\usepackage{hyperref}
\newtheorem{theorem}{Theorem}
\newtheorem{lemma}[theorem]{Lemma}
\newcommand{\C}{\mathbb C}
\newcommand{\R}{\mathbb R}
\newcommand{\Pp}{\mathcal P}
\newcommand{\dd}{\,d}
\title{The finite mixed and confluent Christoffel determinant\newline on the original cyclic source}
\author{}
\date{20 September 2026}
\begin{document}
\maketitle

\section{Status, source, and exact receiving object}\label{sec:status}

This note derives a finite identity from the equal-size case of G. Akemann
and G. Vernizzi, \emph{Characteristic Polynomials of Complex Random Matrix
Models}, arXiv:hep-th/0212051v2, original author source
\texttt{AkemannVernizzi.tex}, labels \texttt{Th}, \texttt{Ddef},
\texttt{proof1}, and \texttt{deg}. Its canonical disk-source identifier is
\texttt{PUBUNIT-688110C14C22CB41B5445BED}. The original source's formula
\texttt{deg} explicitly retains the factorial product for coincident
arguments. We use divided derivatives throughout, so those factorials
occur in the entries rather than as an outside multiplier.

The programme already proved the diagonal source/boundary determinant
quotient in \emph{Source--boundary volume dynamics for the arithmetic
control}, supplied in \texttt{Tau\_Toda\_Volume\_Control\_2026-09-12.zip},
member \texttt{Tau\_Toda\_Volume\_Control/NOTE.tex}, equations (11)--(16).
The present diagonal formula is an exact kernel-jet re-expression of
that quotient. The mixed formula retains two possibly different monic
relations and their complex cross-products. No new general integrable
system theorem, growing-degree estimate, or arithmetic-zero conclusion
is asserted.

Fix the original finite full-order packet \(h\), tensor order \(k\geq1\),
and source
\begin{equation}\label{eq:source}
 c=\frac{k}{2},\qquad S=c+iu,\qquad
 v_h=\frac{2\xi}{h},\qquad
 w_h(t)=\frac{|v_h(1/2+it)|^2}{2\pi},\qquad
 d\mu(u)=m_{h,k}(u)\,du,\qquad m_{h,k}=w_h^{*k}.
\end{equation}
Thus \(\mu(\R)=\mu_h^k\), where \(\mu_h=\int_\R w_h(t)\,dt\).
This measure is never divided by its mass.
The original polynomial source and arithmetic map are
\begin{equation}\label{eq:original-maps}
 \mathcal V_{h,k}P
   =P(D_1+\cdots+D_k)(F_h^{\otimes k}),\quad
 \mathcal M F_h=v_h,\quad
 J^{(k)}\mathcal V_{h,k}=\eta\pi_\chi,
 \quad
 \eta[P]=\upsilon_h^{\otimes k}
             P\!\left(\sum_iM_{s_i}\right)1.
\end{equation}
Here \(\upsilon_h=j_h(v_h)\) is the original invertible Taylor unit,
\(\chi=\chi_{h,k}\) is the actual monic cyclic relation,
\(q=\deg\chi\), \(C=\C[S]/(\chi)\), and \(\eta\) is the original
injection. The norm of this very source is
\begin{equation}\label{eq:inner}
 \langle P,Q\rangle
 =\int_\R\overline{P(c+iu)}Q(c+iu)\,d\mu(u)
 =\langle\mathcal V_{h,k}P,\mathcal V_{h,k}Q\rangle.
\end{equation}
The inner product is conjugate-linear in its first argument.

All polynomial moments of \(\mu\) are finite, and every nonzero
polynomial has positive norm. The following sufficient bound and its
alternating-series/Stirling derivation are those of the original Toda
note, Section 3.1, equations (7)--(8), retained here at their point of
use. For completeness, its bound on the critical line is
\[
 w_h(t)\leq C_h(1+|t|)^{11/2-2\deg h}e^{-\pi|t|/2}
 \quad\text{outside a compact interval}.
\]
Indeed the alternating Dirichlet series has bounded partial sums, so
summation by parts on \(\Re s=1/2\) gives an \(O(1+|s|)\) bound; the
factor \(1-2^{1-s}\) has modulus at least \(\sqrt2-1\); Stirling's
vertical estimate, \(s(s-1)\), and monic \(h\) give precisely the
displayed sufficient bound. On the compact interval the full-order
division makes \(v_h\) entire. It is not the zero function, and its
zeros on the line form a discrete set. Thus \(w_h>0\) almost everywhere.
For \(0<a<\pi/2\), Fubini and \(|t_1+\cdots+t_k|\leq\sum|t_i|\) give
\[
 \int e^{a|u|}\,d\mu(u)
 \leq\left(\int e^{a|t|}w_h(t)\,dt\right)^k<\infty.
\]
The convolution is positive almost everywhere: for each further
convolution the integrand is positive for almost every integration
variable, since translations preserve null sets. The density is finite
almost everywhere by its finite mass. A nonzero polynomial vanishes at
only finitely many points of the line. Its squared norm is therefore
strictly positive. These facts justify all finite integrations below.

\section{Conventions and the finite identity}\label{sec:statement}

Write \(\Pp_j=\C[S]_{\leq j}\), with \(\Pp_{-1}=\{0\}\).
For \(d\geq0\), let
\[
 e_d(S)=(1,S,\ldots,S^{d-1})^T,\qquad
 H_d=\int\overline{e_d(S)}e_d(S)^T\,d\mu,
 \qquad D_d=\det H_d.
\]
An empty vector and a \(0\times0\) determinant have their usual meanings,
and \(D_0=1\). A coefficient column \(a\) represents \(e_d(S)^Ta\),
whose squared norm is \(a^*H_da\). Positivity makes \(H_d\) invertible
for \(d>0\). Gram--Schmidt in the original monomial order produces the
unique monic orthogonal \(p_j\); let \(\omega_j=\|p_j\|^2>0\). Its
unit-triangular coefficient matrix has determinant one. Hence
\begin{equation}\label{eq:kernel}
 D_d=\prod_{j=0}^{d-1}\omega_j,
 \qquad
 K_d(v,\overline w)=\sum_{j=0}^{d-1}
      \frac{p_j(v)\overline{p_j(w)}}{\omega_j}
   =e_d(v)^TH_d^{-1}\overline{e_d(w)}.
\end{equation}
For the final equality, the coefficient matrix \(A\) of the \(p_j\)
satisfies \(A^*H_dA=\operatorname{diag}(\omega_j)\); solving for
\(H_d^{-1}\) and multiplying the evaluation rows proves the formula.
The kernel here is the source's \emph{bare} polynomial kernel, not a
kernel with external density factors attached.

Let \(\psi\) be another monic polynomial of the same degree \(q\).
Factor the actual polynomials over \(\C\), retaining all multiplicities:
\[
 \chi(S)=\prod_{a=1}^{A}(S-v_a)^{m_a},\qquad
 \psi(S)=\prod_{b=1}^{B}(S-w_b)^{\ell_b},\qquad
 \sum_a m_a=\sum_b\ell_b=q.
\]
Roots are distinct within either displayed list; roots between the lists
may coincide. Choose an order of each root list, and order each root's
derivative rows increasingly from zero. Define
\begin{equation}\label{eq:V}
 V_\chi=\prod_{a<b}(v_b-v_a)^{m_am_b},\qquad
 V_\psi=\prod_{a<b}(w_b-w_a)^{\ell_a\ell_b}.
\end{equation}
For \(n\geq0\), put \(d=n+q\) and define the \(q\times q\) matrix
\begin{equation}\label{eq:C}
 C_n(\chi,\psi)_{(a,r),(b,s)}
  =\frac{1}{r!s!}
     \partial_v^r\partial_{\overline w}^{s}
           K_d(v,\overline w)\big|_{v=v_a,w=w_b},
 \quad 0\leq r<m_a,\quad0\leq s<\ell_b.
\end{equation}
The conjugate variable is treated as an independent polynomial variable
when differentiating. Thus its derivative is the complex conjugate of
the corresponding ordinary derivative of \(p_j\).

The requested mixed determinant, with its exact displayed orientation,
is
\begin{equation}\label{eq:B}
 B_n(\chi,\psi)=
 \det\left[\int S^i\chi(S)\overline{S^j\psi(S)}\,d\mu\right]_{i,j=0}^{n-1},
 \qquad B_0=1.
\end{equation}
If \(T_\chi,T_\psi:\C^n\to\C^d\) are the coefficient matrices of
multiplication by their polynomials, the matrix displayed in
\eqref{eq:B} is \((T_\psi^*H_dT_\chi)^T\). In particular,
\begin{equation}\label{eq:B-orientation}
 B_n(\chi,\psi)=\det(T_\psi^*H_dT_\chi),\qquad
 B_n(\psi,\chi)=\overline{B_n(\chi,\psi)}.
\end{equation}
The two cross-Gram orientations must not be identified entry by entry.

\begin{theorem}[Finite mixed/confluent re-expression]\label{thm:mixed}
For the unchanged source \eqref{eq:source} and all the integers and
polynomials just specified,
\begin{equation}\label{eq:main}
 \boxed{B_n(\chi,\psi)
   =D_{n+q}\frac{\det C_n(\chi,\psi)}{V_\chi\overline{V_\psi}}.}
\end{equation}
An empty root list has \(V=1\), and the empty kernel determinant is one.
The mixed determinant can be complex or zero; no positivity of a mixed
matrix is used. For \(\psi=\chi\), its value is strictly positive when
\(n>0\).
\end{theorem}

\section{Finite determinant integration, with signs}\label{sec:integration}

Here and below \(z_t=c+iu_t\) remains a point on the original vertical
line, integrated with \(d\mu(u_t)\). No area measure or matrix ensemble
is introduced. For two ordered polynomial lists \(f_i,g_i\),
\(0\leq i<n\), expansion of the two determinants gives
\begin{equation}\label{eq:andreief}
 \frac1{n!}\int
 \det[f_i(z_t)]_{t,i}\,
 \overline{\det[g_i(z_t)]_{t,i}}
 \prod_{t=1}^n d\mu(u_t)
  =\det\left[\int f_i(S)\overline{g_j(S)}\,d\mu\right]_{i,j}.
\end{equation}
To verify the factor, a pair of permutations \(\sigma,\tau\) contributes
\(\operatorname{sgn}(\sigma)\operatorname{sgn}(\tau)
\prod_t M_{\sigma(t),\tau(t)}\), where
\(M_{ij}=\int f_i\overline{g_j}\,d\mu\).
Reindexing by \(i=\sigma(t)\) makes the second index
\((\tau\sigma^{-1})(i)\). For each relative permutation there are
exactly \(n!\) choices of \(\sigma\), giving \(n!\det M\).
All sums are finite and all products are absolutely integrable.
For \(n=0\), both sides are one.

First assume all roots within each polynomial are simple and list them
as \(v_1,\ldots,v_q\) and \(w_1,\ldots,w_q\). Use
\[
 \Delta_t(x_1,\ldots,x_t)=\det[x_a^{j-1}]_{a,j=1}^t
 =\prod_{a<b}(x_b-x_a).
\]
With the integrated rows first and the external rows last,
\[
 \Delta_d(z,v)=\Delta_n(z)\Delta_q(v)
                    \prod_{t=1}^n\prod_{a=1}^q(v_a-z_t).
\]
Thus
\[
 \det[z_t^i\chi(z_t)]_{t,i}
 =\Delta_n(z)\prod_t\chi(z_t)
 =(-1)^{nq}\frac{\Delta_d(z,v)}{\Delta_q(v)}.
\]
The corresponding conjugate expression has the same real sign
\((-1)^{nq}\). Their product has sign \((-1)^{2nq}=1\). Consequently
\begin{equation}\label{eq:augmented}
 B_n(\chi,\psi)=\frac{1}{n!\Delta_q(v)\overline{\Delta_q(w)}}
  \int\Delta_d(z,v)\overline{\Delta_d(z,w)}\prod_t d\mu(u_t).
\end{equation}

Let \(\phi_j=p_j/\sqrt{\omega_j}\), taking the positive square roots of
the actual norms. Replacing monomials by their monic orthogonal
polynomials has determinant one, so
\[
 \Delta_d(z,v)=D_d^{1/2}\det[\phi_j(x_a)]_{a=1,\ldots,d; j=0,\ldots,d-1},
 \quad x=(z,v).
\]
Expand this determinant along its last \(q\) rows. For a \(q\)-element
column set \(J\), in increasing order, the external minor is
\(\det\Phi_{v,J}\), and the integrated minor uses the complementary
increasing column set \(I\). The Laplace sign is the sign determined by
the sum of the selected row and column indices. Apply
\eqref{eq:andreief} to a pair of integrated minors. Orthonormality gives
\[
 \int\det\Phi_{z,I}\overline{\det\Phi_{z,I'}}\prod_t d\mu(u_t)
   =n!\det[\delta_{ij}]_{i\in I,j\in I'}
   =n!\,\mathbf1_{I=I'}.
\]
When \(I\ne I'\), the indicated square submatrix has a zero row;
when they agree it is the identity. Hence only \(J=J'\) survives, and
its two Laplace signs multiply to one. The result is
\[
 \int\Delta_d(z,v)\overline{\Delta_d(z,w)}\prod_t d\mu(u_t)
  =n!D_d\sum_{|J|=q}\det\Phi_{v,J}\overline{\det\Phi_{w,J}}
  =n!D_d\det(\Phi_v\Phi_w^*).
\]
The last equality follows by expanding the determinant of the product:
terms with a repeated selected column cancel in pairs, and the remaining
terms group into the displayed products of increasing-column minors.
Its entries are \(K_d(v_a,\overline w_b)\). Substituting in
\eqref{eq:augmented} proves \eqref{eq:main} for simple roots. This proof
also applies at \(n=0\), when there are no integrated rows.

\section{Confluence and factorials}\label{sec:confluence}

We prove the exact confluent determinant rule used here. Let a cluster
of \(m\) ordered row arguments tend to \(v\) as
\(v+\varepsilon\alpha_t\), \(t=0,\ldots,m-1\), with distinct constants
\(\alpha_t\). For a polynomial row function \(F\), Taylor expansion is
the finite identity
\[
 F(v+\varepsilon\alpha_t)
  =\sum_{r\geq0}\varepsilon^r\alpha_t^r\frac{F^{(r)}(v)}{r!}.
\]
Multilinearity of the determinant kills every term with repeated
derivative order within the cluster. The smallest possible total order
is \(0+1+\cdots+(m-1)=m(m-1)/2\). At that order, summing over the
permutations of these orders produces exactly
\[
 \varepsilon^{m(m-1)/2}
  \det[\alpha_t^r]_{t,r=0}^{m-1}
  \det\left[\frac{F^{(r)}(v)}{r!}\right]_{r=0}^{m-1}
\]
with the other rows left in their positions. There is no additional
sign: both the rows in the scalar Vandermonde and the derivative rows
have the increasing orders specified above. Higher-order terms carry
strictly larger powers of \(\varepsilon\).

Performing this operation in all clusters makes the within-cluster
factors of the full Vandermonde exactly cancel the displayed scalar
Vandermonde factors and powers of \(\varepsilon\). Its between-cluster
factors tend to \(\prod_{a<b}(v_b-v_a)^{m_am_b}=V_\chi\).
Thus the limiting row determinant is the determinant of divided
derivatives, divided by \(V_\chi\). Applying the identical argument to
columns in the independent variable \(\overline w\) gives the column
factorials and \(\overline{V_\psi}\). The roots within either cluster
can be approached along the stated distinct paths; no real ordering or
separation between the two polynomials is required.

The left side of \eqref{eq:main} is a polynomial in the coefficients of
\(\chi\) and the conjugates of the coefficients of \(\psi\), with
finite moments as coefficients. It therefore tends to the exact
repeated-root determinant. The right side tends, by the preceding
calculation, to \eqref{eq:C} divided by
\(V_\chi\overline{V_\psi}\). This completes the proof of
Theorem~\ref{thm:mixed}.

If ordinary derivatives are used instead, define
\[
 F_\chi=\prod_a\prod_{r=0}^{m_a-1}r!,\qquad
 F_\psi=\prod_b\prod_{s=0}^{\ell_b-1}s!.
\]
The ordinary-derivative matrix has determinant
\(F_\chi F_\psi\det C_n\). Formula \eqref{eq:main} then has the
additional denominator \(F_\chi F_\psi\).
In the source's all-equal, unbarred example \texttt{deg}, this is
\(\prod_{j=1}^{q-1}j!\). Its explicit factorial is consistent with the
present calculation; the brief confluent prose is not being used to
allege an error in the source.

\section{Literal jet and remainder maps}\label{sec:maps}

Set \(N=d-1=n+q-1\). Define the divided-jet evaluation map
\begin{equation}\label{eq:E}
 E_\chi:\Pp_N\longrightarrow\C^q,
 \qquad P\longmapsto
       \left(\frac{P^{(r)}(v_a)}{r!}\right)_{(a,r)}.
\end{equation}
The same symbol denotes its matrix in the original coefficient basis:
\[
 (E_\chi)_{(a,r),j}=
 \begin{cases}\binom jr v_a^{j-r}&j\geq r,\\0&j<r.\end{cases}
\]
It has kernel exactly \(\chi\Pp_{n-1}\). Indeed Taylor's finite
expansion at \(v_a\) shows that all its first \(m_a\) divided jets
vanish if and only if \((S-v_a)^{m_a}\) divides \(P\). Distinct linear
factors are coprime, so their product divides \(P\); division decreases
the degree by exactly \(q\) for a nonzero polynomial. The converse
follows from the same expansion. At \(n=0\) this kernel is zero.

Let \(r_\chi:\Pp_N\to\Pp_{q-1}\) be the unique remainder after monic
division by \(\chi\), and let \(J_\chi\) be the \(q\times q\) matrix
of \(E_\chi\) restricted to \(\Pp_{q-1}\). The preceding divisibility
proof makes this restriction injective. Equal finite dimensions make it
bijective, and
\begin{equation}\label{eq:E-factor}
 E_\chi=J_\chi r_\chi,
 \qquad
 \C[S]/(\chi)\xrightarrow{\ [P]\mapsto(P^{(r)}(v_a)/r!)\ }
 \bigoplus_a\C[\epsilon_a]/(\epsilon_a^{m_a})
\end{equation}
is an algebra isomorphism, not only a vector-space identification.
For multiplication, the coefficient of \(\epsilon_a^r\) in the product
of the two truncated Taylor series is
\(\sum_{t=0}^r P^{(t)}(v_a)Q^{(r-t)}(v_a)/(t!(r-t)!)\), which is
\((PQ)^{(r)}(v_a)/r!\) by the Leibniz formula. This proves the algebra
claim and fixes the target multiplication exactly.

Confluence applied to the ordinary evaluation matrix
\([v_a^j]\) gives
\begin{equation}\label{eq:detJ}
 \det J_\chi=V_\chi.
\end{equation}
This also proves the exact orientation of its determinant. Permuting two
root blocks of sizes \(m_a,m_b\) changes both sides by
\((-1)^{m_am_b}\); changing a root ordering has no effect on
\eqref{eq:main}. Increasing derivative order within each block remains
part of the definition.

Formula \eqref{eq:kernel} and differentiation give the literal matrix
identity
\begin{equation}\label{eq:C-matrix}
 C_n(\chi,\psi)=E_\chi H_d^{-1}E_\psi^*.
\end{equation}
If \(k_{a,r}(S)=\sum_{j<d}p_j(S)
\overline{p_j^{(r)}(v_a)}/(r!\omega_j)\), then
\(\langle k_{a,r},P\rangle=P^{(r)}(v_a)/r!\).
Consequently \(C_n(\chi,\psi)\) is the standard, conjugate-first
cross-Gram of these representing vectors, with the \(\chi\) vector
in its first slot and the \(\psi\) vector in its second slot.
This is consistent with \eqref{eq:B-orientation}: the requested boundary
matrix is transposed when expressed in coefficient cross-Gram notation,
while the evaluation-kernel matrix is not.

\section{The unchanged least-source-norm quotient}\label{sec:quotient}

Now put \(\psi=\chi\), and abbreviate
\(E=E_\chi\), \(J=J_\chi\), \(C_n=C_n(\chi,\chi)\).
The surjectivity of \(E\) and positivity of \(H_d\) imply that \(C_n\)
is positive definite for \(q>0\): for a nonzero column \(y\),
\(E^*y\ne0\), and \(y^*C_ny=(E^*y)^*H_d^{-1}(E^*y)>0\).
For a prescribed jet column \(y\), define
\begin{equation}\label{eq:jet-section}
 a_y=H_d^{-1}E^*C_n^{-1}y.
\end{equation}
Then \(Ea_y=y\). If \(Eb=0\),
\(a_y^*H_db=y^*C_n^{-1}Eb=0\). Every other lift is \(a_y+b\), and
\[
 \|a_y+b\|_{H_d}^2=\|a_y\|_{H_d}^2+\|b\|_{H_d}^2.
\]
Thus \eqref{eq:jet-section} is the unique minimum-norm lift of the
specified jet column; its norm is \(y^*C_n^{-1}y\).

A remainder coefficient column \(x\in\C^q\), in the exact ordered basis
\(1,S,\ldots,S^{q-1}\), has jet column \(Jx\). Its unique
least-source-norm polynomial representative and quotient Gram are
\begin{equation}\label{eq:canonical-section}
 \widehat R_N=H_d^{-1}E^*C_n^{-1}J,\qquad
 G_N=J^*C_n^{-1}J,\qquad R_N=\mathcal V_{h,k}\widehat R_N.
\end{equation}
These are precisely the original quotient maps, because
\(r_\chi\widehat R_N=J^{-1}E\widehat R_N=I\), their images are
orthogonal to the exact kernel \(\chi\Pp_{n-1}\), and the preceding
strict norm decomposition proves uniqueness. In particular the
original arithmetic map is literally retained:
\[
 J^{(k)}R_N=\eta\pi_\chi\widehat R_N=\eta.
\]
No Taylor unit has been replaced by one, and no target jet norm has been
substituted for the source norm. The jet norm \(C_n^{-1}\) was deduced
from that source norm by minimization.

The determinant of \eqref{eq:canonical-section}, followed by
Theorem~\ref{thm:mixed}, is
\begin{equation}\label{eq:quotientdet}
 \boxed{\det G_N=\frac{|V_\chi|^2}{\det C_n(\chi,\chi)}
              =\frac{D_{n+q}}{B_n(\chi,\chi)},
       \qquad N=n+q-1.}
\end{equation}
For an independent comparison with the earlier Toda note, concatenate
the coefficient columns of
\(1,S,\ldots,S^{q-1}\) and
\(\chi,S\chi,\ldots,S^{n-1}\chi\).
Its leading-degree coefficient matrix is unit triangular, of determinant
one. In this basis the source Gram has lower-right block
\(T_\chi^*H_dT_\chi\) of determinant \(B_n(\chi,\chi)\).
Subtracting from every initial lift its orthogonal projection to these
boundary columns is a determinant-one column operation, with the
matching conjugate row operation. It produces the orthogonal block Gram
with diagonal blocks \(G_N\) and \(T_\chi^*H_dT_\chi\).
Therefore \(D_d=\det G_N\,B_n(\chi,\chi)\).
This is the original note's equation (16), now accompanied by its exact
evaluation-jet factorization.

\subsection{The norm of a monic polynomial in the original relation source}
\label{sec:relationnorm}

Equip the original input polynomials to \(M_\chi\) with the pulled-back
inner product
\[
 \langle P,Q\rangle_\chi=\langle\chi P,\chi Q\rangle
 =\int\overline{P(S)}Q(S)|\chi(S)|^2\,d\mu.
\]
This is the squared norm of the actual multiplication map into the same
source, not a replacement of \(d\mu\) by a probability measure. Its
moments are finite and it is positive definite: multiplying a nonzero
polynomial by the nonzero polynomial \(\chi\) cannot yield zero, and
the original norm is positive definite. Let \(p_j^\chi\) be its monic
orthogonal polynomials, with squared norms \(\nu_j\).
The unit-triangular change from \(1,S,\ldots,S^{n-1}\) to
\(p_0^\chi,\ldots,p_{n-1}^\chi\) gives
\[
 B_n(\chi,\chi)=\prod_{j=0}^{n-1}\nu_j.
\]
Dividing this equality at consecutive orders, and then substituting
\eqref{eq:main}, proves for every \(n\geq0\)
\begin{equation}\label{eq:relationnorm}
 \boxed{\nu_n=\frac{B_{n+1}(\chi,\chi)}{B_n(\chi,\chi)}
       =\omega_{n+q}
          \frac{\det C_{n+1}(\chi,\chi)}
               {\det C_n(\chi,\chi)}.}
\end{equation}
The factors \(|V_\chi|^2\) cancel because the relation and its exact
multiplicities are unchanged at consecutive orders. The denominator
is positive; for \(q=0\) both kernel determinants are empty and
\(\nu_n=\omega_n\). Under a scalar comparison
\(d\mu\mapsto\lambda\,d\mu\), the kernel-determinant ratio is
unchanged and \(\omega_{n+q}\) scales by \(\lambda\), so the full
original mass is retained in \(\nu_n\).

Akemann--Vernizzi's source discusses the general-size expression at
\texttt{Lnorm} as expected and checked in several small sizes; the
equal-degree determinant theorem above plus this finite Gram argument
supplies its norm expression for the present original relation source,
including repeated roots. This is not a priority claim about the later
literature, and it does not claim that every separate orthogonal-polynomial
formula discussed there has been proved here. The original Toda note
already identifies \(\nu_n=B_{n+1}/B_n\) in its equation (19).

\section{Boundary cases}\label{sec:boundary}

If \(n=0\) and \(q>0\), then \(d=q\), \(E_\chi=J_\chi\), and
\[
 C_0(\chi,\psi)=J_\chi H_q^{-1}J_\psi^*,\qquad
 \det C_0=\frac{V_\chi\overline{V_\psi}}{D_q}.
\]
Thus \eqref{eq:main} reads \(1=1\), not an unstated limiting convention.
There are no admitted boundary columns, the quotient section is the
identity in coefficients, and \(G_{q-1}=H_q\).

If \(q=0\), both monic degree-zero polynomials are \(1\), their root
lists and jets are empty, and \(B_n(1,1)=D_n\). Both sides of
\eqref{eq:main} agree. The quotient \(\C[S]/(1)\) is the zero module,
the jet target is the zero vector space, and its determinant is one.
Its least-norm representative of its sole class is zero; all of
\(\Pp_{n-1}\) is the relation subspace. No positive-dimensional inverse
is asserted. For \(n=q=0\), take \(\Pp_{-1}=0\); every determinant in
the statement equals one. One may read all \(0\times0\) maps as the
unique maps on the zero space, but no inverse formula is needed to prove
this case.

\section{Mass scaling without changing the source}\label{sec:mass}

To expose every mass power, consider only as an algebraic comparison
\(d\mu\mapsto\lambda\,d\mu\), with \(\lambda>0\). Monic
orthogonal polynomials stay the same, their norms become
\(\lambda\omega_j\), and
\[
 H_d\mapsto\lambda H_d,\quad D_d\mapsto\lambda^dD_d,
 \quad K_d\mapsto\lambda^{-1}K_d,\quad
 C_n\mapsto\lambda^{-1}C_n,\quad B_n\mapsto\lambda^nB_n.
\]
The right side of \eqref{eq:main} scales as
\(\lambda^{d-q}=\lambda^n\), exactly as its left side.
The quotient section \eqref{eq:canonical-section} stays unchanged and
\(G_N\mapsto\lambda G_N\), so its determinant scales as
\(\lambda^q\). These identities show where the original total mass
\(\mu_h^k\) is retained. No such scaling is actually performed on
\eqref{eq:source}.

\section{The exact \(S=c+iu\) coordinate maps and phases}\label{sec:phases}

The coefficient change from \(S\) to \(u\) is the upper triangular
matrix
\[
 (L_d)_{rj}=\begin{cases}\binom jr c^{j-r}i^r&r\leq j,\\0&r>j,
 \end{cases}
 \qquad 0\leq r,j<d,
 \qquad\det L_d=i^{d(d-1)/2}.
\]
It is defined by the exact equality
\(P(c+iu)=e_d(u)^TL_da\) for \(P(S)=e_d(S)^Ta\).
If \(H_d^{u}\) is the original \(u\)-monomial Gram with the same
\(d\mu(u)\), then
\begin{equation}\label{eq:coordinategram}
 H_d=L_d^*H_d^{u}L_d,\qquad D_d=\det H_d^{u}.
\end{equation}
The determinant equality follows only after retaining and multiplying
the phases \(\overline{\det L_d}\det L_d=1\).

Let \(Q_j\) be the monic \(u\)-orthogonal polynomial. The polynomial
\(i^jQ_j((S-c)/i)\) is monic in \(S\), has the required
orthogonality, and has norm \(\|Q_j\|^2\); uniqueness gives
\[
 p_j(S)=i^jQ_j((S-c)/i),\qquad \omega_j=\|Q_j\|^2.
\]
For \(\zeta=(v-c)/i\), \(\eta=(w-c)/i\), each summand of the kernel
has the paired phases \(i^j\overline{i^j}=1\). Hence
\[
 K_d(v,\overline w)=K_d^u(\zeta,\overline\eta),\qquad
 \partial_v=i^{-1}\partial_\zeta,\quad
 \partial_{\overline w}=\overline{i}^{\,-1}\partial_{\overline\eta}
                         =i\partial_{\overline\eta}.
\]
Define the chart polynomials
\(\widetilde\chi(u)=i^{-q}\chi(c+iu)\) and
\(\widetilde\psi(u)=i^{-q}\psi(c+iu)\). They are monic; the actual
source multipliers remain
\(\chi(c+iu)=i^q\widetilde\chi(u)\) and
\(\psi(c+iu)=i^q\widetilde\psi(u)\).
For their roots \(\zeta_a=(v_a-c)/i\) and
\(\eta_b=(w_b-c)/i\), put
\[
 R_\chi=\sum_a\frac{m_a(m_a-1)}2,\quad
 P_\chi=\sum_{a<b}m_am_b,\quad
 R_\chi+P_\chi=\frac{q(q-1)}2,
\]
and similarly for \(\psi\). Let \(\Lambda_\chi\) be the diagonal
jet matrix with entries \(i^{-r}\) in row \((a,r)\).
The chain rule gives exact matrix identities
\begin{equation}\label{eq:jetphases}
 E_\chi^S=\Lambda_\chi E_{\widetilde\chi}^{u}L_d,
 \quad J_\chi^S=\Lambda_\chi J_{\widetilde\chi}^{u}L_q,
 \quad C_n^S(\chi,\psi)
       =\Lambda_\chi C_n^u(\widetilde\chi,\widetilde\psi)
                          \Lambda_\psi^*.
\end{equation}
Their determinant phases are
\[
 \det C_n^S=i^{-R_\chi+R_\psi}\det C_n^u,\qquad
 V_\chi^S=i^{P_\chi}V_{\widetilde\chi}^{u},\qquad
 \overline{V_\psi^S}=i^{-P_\psi}\overline{V_{\widetilde\psi}^{u}}.
\]
The resulting phase of the ratio in \eqref{eq:main} is
\(i^{-R_\chi+R_\psi-P_\chi+P_\psi}=1\), because both polynomials
have the same degree \(q\). This checks the mixed orientation even
when their multiplicity patterns differ.

For the boundary determinants, changing the \(n\) input monomials
introduces \(L_n\) and its adjoint; the two degree-\(q\) multiplier
phases contribute \(\overline{i^q}i^q=1\) to the coefficient
cross-Gram. Thus
\[
 B_n^S(\chi,\psi)=B_n^u(\widetilde\chi,\widetilde\psi)
\]
with the measure still \(d\mu\). The quotient map is the literal
algebra isomorphism
\[
 \C[S]/(\chi)\longrightarrow\C[u]/(\widetilde\chi),
 \qquad[P(S)]\longmapsto[P(c+iu)].
\]
It is well defined since the relation becomes
\(i^q\widetilde\chi\); its inverse substitutes \(u=(S-c)/i\).
It preserves source norms by \eqref{eq:inner}. On remainder coefficient
columns it is \(L_q\), so
\(G_N^S=L_q^*G_N^uL_q\), again with determinant multiplier one.
This is a proved coordinate transport, not a new relation or a new
measure.

\section{What has and has not been checked}\label{sec:checks}

The companion \texttt{check\_identities.py} uses exact rational complex
arithmetic to verify determinant integration, mixed orientations,
confluent derivatives, nonunit mass powers, all boundary cases, and the
minimum-section equations. Its finite atomic fixtures test the
polynomial and matrix identities proved above; they are not substituted
for the programme density, and they are not a numerical computation of
the original arithmetic moments. The proof applies to those original
moments by Sections~\ref{sec:status}--\ref{sec:integration}.
No probabilistic ensemble interpretation is required at any stage.

\end{document}
